%%%%%%%%%%%%%%%%%%%%%%%%%%%%%%%%%%%%%%%%%%%%%%%%%%%%%%%%%%%%%%%%%%%%%%%%%%%%%%%%
%%
\cleardoubleoddpage%  Make sure to start each chapter on a new odd page
\chapter{Method}
\label{ch:method}


\section{Overview: One Framework, One Swappable Slot}
\label{sec:method:overview}

% TODO: one pipeline figure with the operator slot drawn as the only part
% that varies: state in -> (fixed) conditioning -> f_theta -> (fixed)
% Delta-t-scaled residual update -> next state. State the design contract:
% loss, curriculum, evaluation, model selection and bookkeeping live in one
% shared pipeline that a backbone may not duplicate. Point to every
% subsection below.


\section{The Fixed Framework}
\label{sec:method:framework}

% TODO: what every arm shares. One visit in, no history (and why).

\subsection{\texorpdfstring{$\Delta t$}{Delta t}-Conditioned Residual Formulation}
\label{sec:method:dt-residual}

% TODO: u_{t+Delta t} = u_t + Delta t * f_theta(u_t, Delta t); why this
% replaces absolute time; the Delta t -> 0 persistence limit; the
% zero-initialised head, so every model starts at exact persistence.

\subsection{Multi-Channel State Propagation}
\label{sec:method:multichannel}

% TODO: how the C = 11 channels move through the pipeline at the tensor
% level. Mathematics and tensors only -- the clinical meaning of the
% channels is established in Section~\ref{sec:data:state}.


\section{Backbone I: The GA-Adapted Graph Solver}
\label{sec:method:mppde}

% TODO: the arm the project built rather than imported: encoder, message
% passing, decoder, Delta-t-conditioned output. Temporal bundling removed
% entirely (strict K = 1).

\subsection{Architecture}
\label{sec:method:mppde:arch}

% TODO: encoder, message and update functions, conditioning vector,
% normalisation, decoder.

\subsection{Graph Construction on an Anisotropic Grid}
\label{sec:method:graph}

% TODO: the load-bearing part of the section. Index-space k-NN (v1) versus
% the dilated physical-scale stencil (v2 lock); offsets derived from the
% lesion scale.

\subsection{Differences from the Original MP-PDE}
\label{sec:method:mppde:diff}

% TODO: summary table of every departure from Brandstetter et al. (2022).


\section{Backbone II: The Countermodel Family}
\label{sec:method:family}

% TODO: the shared backbone contract; U-Net at two capacities, FNO, FNO +
% 3x3 local kernel; graph U-Net; per-pixel floor; Finite Element Network,
% free-form and with the transport term; the Runge--Kutta wrapper;
% parameter matching (1.25x band, and the two deliberate out-of-band
% controls). Same depth as Backbone I.


\section{The Moving-Mesh Extension as a Tested Hypothesis}
\label{sec:method:mmpde}

% TODO: DMM + dual branch, framed as a hypothesis the design was built to
% measure. The alpha gate (zero-initialised) is the instrument; alpha -> 0
% was a pre-registered reportable outcome. res_cut is NOT part of the
% framework and must not be described as a component.

\subsection{Data-Free Mesh Mover for GA}
\label{sec:method:dmm}

% TODO: physics loss (Monge--Amp\`ere + boundary + convexity), sampling,
% monitor function. The monitor is a plain scalar monitor on the Gaussian-
% blurred mask; the DMM trains and is applied on the native (49, 1024) grid.
% The Frobenius-norm multi-channel monitor and the 256^2 SLO path are
% superseded -- historical steps only, never the design. Pretrained, frozen.

\subsection{Dual-Branch Composition}
\label{sec:method:dual}

% TODO: uniform branch returns the full next state; moved branch returns a
% pure Delta-t delta; ItpNet interpolation renormalised to sum to one; the
% learnable scalar gate alpha.


\section{Conditioning: Covariates and the Coefficient Surrogate}
\label{sec:method:covariates}

% TODO: age and sex broadcast per node into the conditioning vector. Their
% value is an open question at this point; Chapter 5 reports it.

\subsection{Surrogate Equation Encoder}
\label{sec:method:encoder}

% TODO: the LayerEncoder CNN, framed as a learned stand-in for theta_PDE.
% Its stride schedule corrects only ~2x of the ~21:1 anisotropy.


\section{Training}
\label{sec:method:training}

% TODO: channel-weighted MSE + soft-Dice objective; time-budgeted
% pushforward curriculum (unroll depth in elapsed days); ItpNet pretraining
% at epoch 0 for the dual branch; per-module gradient clipping; optimiser
% and schedule. Monotonic penalty is 0.0 and the residual is signed.


\section{Implementation Details}
\label{sec:method:implementation}

% TODO: hardware, framework versions, run identity / checkpointing,
% graph precomputation, notable engineering choices.


%%
%%%%%%%%%%%%%%%%%%%%%%%%%%%%%%%%%%%%%%%%%%%%%%%%%%%%%%%%%%%%%%%%%%%%%%%%%%%%%%%%
